% 本文件由 LaTeX 排版 Agent 根据有效 translation.json 自动生成。
% author=Augustine of Hippo; work_id=1309; title=论婚姻之善

\chapter{论婚姻之善}

% structure: p0001 source=h1 target=omitted reason=page-title-already-rendered-by-parent

这篇论文以及下一篇是针对约维尼安异端尚存的部分而写。圣奥斯定在\WorkTitle{论婚姻与欲望}卷二第二十三章中提到这一错误。他说：\Quote{约维尼安，}他说，\Quote{他几年前试图建立一个新异端，说天主教徒袒护摩尼教徒，因为他反对他们，将神圣的童贞置于婚姻之上。}在他的\WorkTitle{论异端}第八十二章中：\Quote{那异端起源于一位名叫约维尼安的隐修士，在我们这个时代，当我们还年轻的时候。}他补充说，它很快被压制和消灭，大约在公元390年，先在\Place{罗马}受谴责，然后在\Place{米兰}。有教宗西利修就此主题致\Place{米兰}教会的信，以及\Place{米兰}教务会议（由圣盎博罗削主持）给他的答复。热罗尼莫曾驳斥约维尼安，但据说他试图为童贞状态的卓越辩护，却以谴责婚姻为代价。为使奥斯定不受这样的抱怨或诽谤，在谈论童贞的优越性之前，他认为最好先写一篇论婚姻之善的论文。\par

我们知道这部作品大约完成于401年，这不仅是根据他的\WorkTitle{更正}的次序，也是根据他大约在那年开始的\WorkTitle{创世纪字义}诸卷。因为在\WorkTitle{创世纪字义}第九卷第七章中，他称赞婚姻之善，他说：\Quote{现在这善有三重：忠信、子女和圣事。就忠信而言，要遵守不在婚约之外与别的男人或女人同寝；就子女而言，要慈爱地迎接、友善地养育、虔诚地教导；就圣事而言，婚姻不可离异，被休的男女即使为了子女也不可再与他人结合。这可以说是婚姻的法则，藉此法则，生育得以端庄，或不节制的乖谬得以归于秩序。关于这些，既然我们已在最近出版的\WorkTitle{论婚姻之善}一书中充分讨论，并在那里按照各等级的价值区分了寡妇的节德和童贞的卓越，我们的笔就不再赘述。}这部作品也在\WorkTitle{论功过与罪之赦}第一卷第二十九章中提及。— \Emph{本笃版编者}\par

1. \Emph{凡}作为人类种族的一部分，人性是社会的，并且拥有友谊的能力作为一项伟大而自然的善；为此，天主愿意从一人创造所有的人，使他们不仅因种类的相似，也因亲缘的纽带而维系在其社会中。因此，人类社会最初的天然纽带是男人和妻子。天主并没有各自创造他们，然后像出生时异族那样将他们结合；而是从另一方创造了另一方，在女人被取出的肋骨上为结合的力量设立了一个标记。因为他们并肩相连，一同行走，一同看向他们所走的方向。接着是子女中的团结合一，这是唯一值得的果实，不是男女的结合，而是性交的果实。因为即使在无性交的情况下，也可能在男女之间存在一种友好而真实的统治与服从的合一。\par

2. 现今也无需探究并发表明确意见，论及天主所降福的第一批人的后裔如何能存在，祂说：\BibleQuote{你们要增长，繁殖，充满大地；}倘若他们没有犯罪，而他们的身体因犯罪应得了死亡的状态，且只有可朽的身体才能有性交。因为对此问题已有数种不同意见；若我们必须考察哪一种更符合圣经真理，则需长篇讨论。因此，若他们没有犯罪，是否会无性交而藉全能造物主的恩赐，以其他方式生子？祂能无父母而创造他们本身，能在童贞母腹中形成基督的肉身，且（甚至对不信者而言）能使蜜蜂无性交而繁殖？或者，许多事物是以奥秘和预像的方式讲述，而\BibleQuote{充满大地，治理大地}应按另一种意义理解，即藉生命的圆满和德能得以实现，以致被称为\BibleQuote{增长，繁殖}的本身增长与增多，应理解为心智的进步与德能的丰盈，正如圣咏所载：\BibleBook{圣咏}{}\BibleQuote{祢要以德能在我灵魂中增加我；}而子孙的延续，仅在因罪而将来有死亡之后才赐予人？或者，这些人的身体起初并非属神的，而是属生灵的，以便藉服从的功德后来成为属神的，从而获得不朽，并非在死亡之后——死亡因魔鬼的嫉妒进入世界，成为罪的惩罚——而是在那变化之后，宗徒指明说：\BibleQuote{然后我们这些活着还存留的人，将同他们一起被提到云中，到空中迎接基督；}从而使我们明白，第一对夫妇的身体在最初形成时是可朽的，然而若非犯罪，他们本不会死，正如天主所警告的：如同祂威胁创伤，因为身体可能受伤，但除非做了祂所禁止的事，创伤本不会发生。因此，即使通过性交，也可能产生这样的身体，它们会增长到一定程度，却不会进入老年；甚至进入老年，却不会死亡，直到大地因祝福的繁殖而充满。因为如果天主使以色列人的衣服在四十年中保持完好无损，那么对于服从祂命令者的身体，祂岂不更会赐予某种极其幸福的稳定状态，直到他们被改变为更好的状态，并非通过人的死亡——灵魂离弃身体——而是通过从可朽到不朽、从属生灵到属神的品质的蒙福改变。这些意见中哪些为真，或者是否可以从这些话语中得出其他一些意见，乃是需要长期探究和讨论的事。\par

3. 现在我们这样说：根据我们所知并在其中受造的出生与死亡的状况，男女婚姻是某种善；圣经如此称赞其盟约，以至于妻子被丈夫休弃后，只要丈夫活着，就不许她再婚；丈夫被妻子休弃后，除非与他分离的妻子已死，就不许他另娶。因此，关于婚姻之善——主也在福音中予以确认，不仅因祂禁止休妻除非为奸淫的缘故，也因祂应邀参加婚筵——我们有充分理由探究它为何是善。在我看来，这不只是因为生育子女，也因性别差异本身的自然结合。否则，在老年人中，特别是若他们已失去子女或从未生育，就不应再称为婚姻。然而，在良好的、虽已年迈的婚姻中，尽管男女之间的壮年激情已消退，夫妻之间的爱德秩序却依然充满活力：因为他们越善，就越早藉双方合意开始节制性交；这不是说以后不能做他们愿意做的事，而是称赞他们起初就不愿做他们有能力做的事。因此，若保持各性别之间应有的尊敬与相互服务的信德，虽然双方的肢体已衰颓近乎尸体，但已恰当结合的灵魂的贞洁仍持续，其纯洁愈经考验而愈显，其安稳愈平静而愈安全。婚姻还有此善：即肉体的或青年的无节制，虽有过失，却被引导至生育子女的正用，从而使婚姻结合从情欲的恶中生出某种善。其次，肉体的情欲受到抑制，并在亲子之情的调节下以较为适度的方式冲动。因为当夫妻相合时，若他们意识到自己是父亲和母亲，那么某种庄重便介入其中，使炽热的快感趋于节制。\par

4. 此外，还有这一点：在夫妻彼此偿还婚姻债时，即使他们以过于放纵和不节制的方式提出要求，但他们仍彼此负有同样的忠诚。宗徒赋予这忠诚如此大的权力，以致称之为\Quote{权力}，说：\BibleQuote{女人对自己的身体没有权力，男人却有权；同样，男人对自己的身体也没有权力，女人却有权。}但违反这忠诚就被称为奸淫，即当一方因自己的情欲煽动，或同意他人的情欲，而与另一方违背婚约发生性行为时；如此便破坏了忠诚，而忠诚即使在关乎身体和卑微的事物中，也是灵魂的一大善；因此它无疑应当被置于甚至包含我们生命的身体健康之上。因为，虽然一点糠秕与许多金子相比几乎等于无，但忠诚在糠秕的事情上如金子一样被保持纯洁时，并不因为它被保持在较小的事情上而减少。但当忠诚被用于犯罪时，我们若仍称它为忠诚，那就奇怪了；然而无论它是何种忠诚，如果行为也违背了它，那么行为就更糟糕；除非它因此被放弃，以便能回归真实合法的忠诚，即通过纠正意志的乖僻来改正罪恶。例如，若有人独自不能抢劫另一个人，便找一个同伴共同行恶，与他约定一起抢劫并分赃；犯罪后，却将全部赃物据为己有。那同伴悲伤抱怨说没有对他保持忠诚，但在他抱怨时，他应当考虑，他自己更应当对人类社会的良好生活保持忠诚，而不应不义地抢劫他人，如果他感受到在罪恶的同伴关系中，对自己未能保持忠诚是多么不义的话。前者在两方面都不忠诚，无疑应被判定更为邪恶。但是，如果他对自己所做的恶事感到不满，并因此不愿与犯罪同伴分赃，以便将赃物归还原主，那么即使是不忠诚的人也不会称他不忠诚。同样，一个女人，如果破坏了婚姻忠诚，却对她的奸夫保持忠诚，当然是邪恶的；但如果连对她的奸夫也不忠诚，那就更坏。此外，如果她悔恨自己的罪，回归婚姻的贞洁，放弃所有奸淫的约定和决议，那么即使奸夫本人都不会认为她是背信弃义之人，我认为这是奇怪的。\par

5. 还有一个常见的问题：当一个男人和一个女人，双方都不是他人的丈夫或妻子，结合并非为了生育，而是由于不节制，只为性交，并且他们之间有这样一种忠诚，即他不与其他女人行此，她也不与其他男人行此，这是否可称为婚姻？也许这不无理由可称为婚姻，如果双方都决定持续到一方死亡，并且虽然他们不是为了生育而结合，但他们也不回避生育，既不愿意无子女，也不用某种恶行手段阻止生育。但如果这两点中有一点或两点都欠缺，我就找不到称它为婚姻的理由。因为，如果一个人暂时娶一个女子，直到他找到另一位配得上他的荣誉或财富、与他地位相当的人结婚；在他的灵魂中他就是一个奸夫，而且不是对他想找到的那位，而是对他与之同床却未与她建立夫妻关系的那位。因此，那女子明知并愿意如此，她与他发生性关系，而与他没有妻子的契约，当然行为不贞。然而，如果她对他保持床笫的忠诚，并且在他结婚后，她本人不再考虑结婚，并准备完全禁绝此类行为，也许我不敢轻易称她为奸妇；但是，既然她知道她与一个不是她丈夫的男人发生性关系，谁能说她无罪呢？此外，如果在那性关系中，就她而言，她只希望生孩子，并且不情愿地忍受超出生育原因的一切，那么有许多已婚妇女比不上她；那些妇女虽然不是奸妇，却强迫她们的丈夫（他们大多也希望操练节制）行肉体之事，不是出于对子女的渴望，而是出于情欲的热焰，无节制地使用她们的权利；但在她们的婚姻中，她们已婚这一点本身就是善。因为她们结婚的目的正是为了使情欲受到合法束缚，不至于无定形、无约束地飘荡；情欲本身有无法克制的肉体软弱，但婚姻有不可解散的忠诚伙伴关系；情欲本身有过度性交的侵犯，但婚姻有贞洁生育的方式。因为，虽然想利用丈夫满足情欲是可耻的，但只愿与丈夫行房、只从丈夫生子却是可敬的。也有男人如此不节制，甚至在妻子怀孕时也不放过她们。因此，已婚夫妇之间任何不端庄、无耻、卑鄙的行为，都是人的罪，而不是婚姻的过错。\par

6. 此外，在更无节制地要求肉身的债务（宗徒不是以命令，而是以允许的方式予以许可）的情况下，即便他们行房事并非为了生育子女，尽管恶习驱使他们如此行，然而婚姻保护他们免于奸淫或私通。因为那并非因婚姻而犯，而是因婚姻而被宽恕。因此，已婚者彼此不仅应在性行为本身中信守忠信，以生育子女——这是人类在这必朽状态中的最初结合；而且，在某种意义上，还应当彼此服侍，扶持彼此的软弱，以避免非法的结合：所以，即使其中一人喜欢永久节欲，也不得未经对方同意而独身。因为至此程度，经上记载：\BibleQuote{妻子没有权力主张自己的身体，而是丈夫；同样，丈夫也没有权力主张自己的身体，而是妻子。}同样，那并非为生育子女，而是因软弱和不节制，不论丈夫向妻子或妻子向丈夫所要求的，他们都不应拒绝，以免因此而陷入该受罚的诱惑，因撒殚的试探，出于双方或其中一方的无节制。因为为生育而行房事是无罪的；但为满足情欲而行，却限于夫妻之间，因床第之忠信，有小罪；但奸淫或私通则有致死之罪，因此，完全节欲甚至比为了生育而行的婚姻房事更好。但因为节欲有更大的功绩，而履行婚姻债务并非罪行，但超出生育需要的索求则是小罪，而行奸淫或私通则是应受惩罚的罪行；已婚者的爱德应当警惕，以免在为自己寻求更大荣誉的机会时，却为自己的伴侣做了导致定罪的事。\BibleQuote{凡休妻的，除非因奸淫的缘故，就是使她犯奸淫。}婚姻契约是如此具有某种圣事性，以至于即使分离也不能使之无效，因为只要她的丈夫还活着，即使被他遗弃，她若另嫁他人就是犯奸淫；而休妻的人正是这恶事的起因。\par

7. 但我感到惊奇，如果像可以休弃犯奸淫的妻子一样，是否也可以休妻后再娶他人。因为圣经在这方面留下了一个难题，因为宗徒说，按主的命令，妻子不应离开丈夫，但若她离开了，应该独身，或与丈夫和好；然而，她显然只应离开犯奸淫的丈夫而独身，以免因离开非犯奸淫的丈夫而使他犯奸淫。但或许她可以公正地与丈夫和好，要么是因他需要被容忍（如果她无法自制），要么是他已经改正。但我不明白，如果男人可以休弃犯奸淫的妻子后再娶，那么女人在离开犯奸淫的丈夫后却不能再嫁。既然这样，婚姻中结合的纽带是如此坚强，尽管它为了生育子女而缔结，但甚至为了生育子女也不能解开。因为男人有权休弃不育的妻子，另娶一个能生育的。然而，这是不允许的；而且在我们这个时代，按照罗马的习惯，也不得再娶另一个妻子，即不能同时拥有多个妻子。而且，如果离开犯奸淫的妻子或丈夫，按理本可以有更多孩子出生，要么女人再嫁，要么男人另娶。然而，如果这不合乎法律，正如神圣的规则似乎规定的那样，有谁不会对此感到惊异，而要学习这婚姻纽带如此强大的意义何在？我绝不认为它具有如此大的效力，除非是从人类这软弱必朽的状态中，取用了某种更大事物的圣事，以致当人离弃它并试图解除它时，它仍保持不散，作为对他们的惩罚。因为婚姻契约并不因离婚而终止；他们即使分离后，仍互为配偶；并且，他们若与别人结合，即使是在自己离婚之后，无论是女人与男人还是男人与女人，都是犯奸淫。然而，在我们天主的圣城，在祂的圣山上，对妻子来说情况并非如此。但是，外邦人的法律不同，谁不知道呢？在那里，通过离婚的介入，在人认知的罪责之外，女人可以嫁给她愿意的人，男人也可以娶他所愿意的人。梅瑟似乎因以色列人的心硬而允许了类似的风俗，即休书之事。在这方面，离婚似乎更多地受到责备，而不是认可。\par

8. \Quote{婚姻应受尊重，}因此，\Quote{床第应是无玷的。} 我们并不将此称为一种善，如同它相对于邪淫是一种善；否则，便会有两种恶，其中第二种更坏；或者邪淫也会成为一种善，因为奸淫更坏；因为侵犯他人的婚姻比与娼妓结合更坏；而奸淫也会成为一种善，因为乱伦更坏；因为与母亲同寝比与别人的妻子同寝更坏；直到我们到达那些事，正如宗徒所说，\Quote{连提及也是可耻的，}所有事相对于更坏的事都将成为善。但谁能怀疑这是错误的呢？因此，婚姻和邪淫并非两种恶，其中后者更坏；而是婚姻和节欲是两种善，其中后者更好，正如暂时的健康和疾病并非两种恶，其中后者更坏；而是健康和永生是两种善，其中后者更好。同样，知识和虚妄并非两种恶，其中虚妄更坏；而是知识和爱德是两种善，其中爱德更好。因为宗徒说：\Quote{知识必要消失，}但在现世它是必要的；而\Quote{爱永存不朽。}同样，这有死的生育——婚姻因此发生——也将消失；但完全脱离一切性行为既是此世天使般的操练，也永远持续。然而，正如义人的宴席胜过亵圣者的禁食，同样，信徒的婚姻应置于不虔敬者的童贞之前。但在这两种情形中，并非宴席优于禁食，而是义德优于亵圣；在此也非婚姻优于童贞，而是信德优于不虔敬。因为义人按需要进食，是为了作为善主给他们的奴隶——\Emph{即，}他们的身体——公平正义之物；但亵圣者禁食是为了服侍魔鬼。同样，信徒结婚是为了贞洁地结合于丈夫，但不虔敬者守童贞是为了远离真天主而行邪淫。因此，正如玛尔大所做的是善事，她忙于服侍圣徒；但她的妹妹玛利亚坐在主的脚前听祂的话，那更是善事；同样，我们赞美苏撒纳在婚姻中的贞洁，但仍将寡妇亚纳的善行置于她之上，更不用说童贞玛利亚。那些用自己的财物供给基督和祂门徒所需的人是行善；但那些舍弃一切财物以便更自由地跟随同一位主的人则行更大的善。但在这些善事中，无论是前者所做的，还是玛尔大和玛利亚所做的，除非另一件事被越过或放弃，更好的事便不能做成。因此我们应当明白，我们不可因此认为婚姻是恶，因为除非从婚姻中戒除，便不能有寡妇的贞洁或童贞的纯洁。因为玛尔大所做的也并非恶，除非她的妹妹戒除它，她便不能做更好的事；同样，接待一位义人或先知到自己家中也不是恶，因为愿意完美跟随基督的人甚至不应有房屋，以便做更好的事。\par

9. 的确，我们必须考虑：天主赐给我们一些善，是因其本身而追求的，如智慧、健康、友谊；另一些则是为其他目的而必需的，如学习、肉食、饮料、睡眠、婚姻、性行为。因为这些中，有些是为智慧而必需的，如学习；有些是为健康而必需的，如肉食、饮料和睡眠；有些是为友谊而必需的，如婚姻或性行为；因为人类之繁衍由此而延续，在其中，友谊的共融是一大善。因此，这些为其他目的而必需的善，谁若不以此为目的使用它们——它们为此而被设立——便是犯罪；有些是小罪，有些是该死的罪。但谁若以此为目的使用它们——它们为此而被赐予——便是行善。因此，对那些不需要它们的人，若他们不使用，便做得更好。所以，当我们有需要时，我们行善而盼望这些善；但我们不盼望比盼望更好：因为当我们视它们为不必要时，我们自身处于更好的状态。因此，结婚是好的，因为生育子女、成为一家之母是好的；但不结婚更好，因为为了人类共融本身而不需要此工是更好的。因为人类种族如今的状态是（那些不能自制的人，不仅忙于婚姻，而且许多人通过非法的姘居而放纵，善造物主从他们的恶中生出善来），从不缺乏众多的后代和丰富的延续，从中可获得神圣的友谊。由此我们得知，在人类最初的时代，主要为了天主子民的繁衍——通过这子民，万民的元首和救主应被预言的并诞生——圣徒们有责任使用这婚姻之善，并非因其本身而追求，而是为其他目的所必需；但如今，为了进入神圣纯洁的共融，从万邦各处有满溢的属灵亲属，甚至那些仅为子女而愿结婚的人也应受劝诫，要他们宁愿使用节欲这更大的善。\par

10. 但我晓得有些人在嘀咕：他们说，如果所有人都禁绝一切性交，人类将从何而来？但愿所有人都这样做，只要\BibleQuote{爱德出于纯洁的心、良好的良心和真诚的信德}；这样，天主的城将更快被充满，世界末日也将提前。因为宗徒在论及此事时，正如明显可见的，他劝勉什么呢？他说：\BibleQuote{我本来愿意众人都像我一样}；或者在那段经文中：\BibleQuote{但是弟兄们，我告诉你们：时间短促了；今后有妻子的，要像没有一样；哭泣的，要像不哭泣的；欢乐的，要像不欢乐的；购买的，要像一无所得的；享用这世界的，要像不享用的。因为这世界的局势正在消逝。我愿意你们无所挂虑。}然后他接着说：\BibleQuote{没有妻子的，挂虑主的事，想怎样悦乐主；娶了妻子的，挂虑世俗的事，想怎样悦乐妻子。未结婚的女子和童贞女不同：未结婚的挂虑主的事，一心使身体和灵魂都圣洁；已婚的则挂虑世俗的事，想怎样悦乐丈夫。}由此在我看来，在这时期，只有那些不能自制的才应当结婚，正如同一宗徒的那句话：\BibleQuote{但如果他们不能自制，就让他们结婚，因为结婚比欲火攻心更好。}\par

11. 然而，婚姻对这些（不能自制）人本身并不是罪；如果选择婚姻是与淫乱相比，那么它将是比淫乱更轻的罪，但仍然会是罪。但现在，我们该如何反驳宗徒最明白的话呢？他说：\BibleQuote{让她随心所欲，不算犯罪，只要她嫁了人；} 以及：\BibleQuote{如果你娶了妻子，不算犯罪；如果童贞女嫁了人，也不算犯罪。} 因此现在确实不可怀疑婚姻不是罪。所以宗徒并不允许婚姻作为\Quote{宽恕的事}；因为谁能否认，说那些获得宽恕的人没有犯罪，是极其荒谬的呢？但他允许作为\Quote{宽恕的事}的是那种因无节制而发生的性交，不仅为了生育子女，有时根本不为了生育子女；并且不是婚姻迫使这种事发生，而是它为此获得了宽恕；然而，只要它没有过分到妨碍那些应当留出来作为祈祷时节的时段，也没有变成那种反乎本性的使用，关于这一点，宗徒在论及不洁和不敬虔之人的过度败坏时，不能沉默。因为为生育所需的性交是无可指责的，并且只有它才配得上婚姻。但那种超出这种必要性的行为，不再遵从理性，而是情欲。然而，婚姻的性质不是要求这个，而是向配偶让步，以免对方因淫乱而犯下该死的罪。但如果双方都被此种情欲支配，他们所做的事显然不是婚姻之事。然而，如果在他们的性交中，他们喜爱正直之事胜过不正直之事，即喜爱属于婚姻之事胜过不属于婚姻之事，那么这是宗徒的权威允许他们作为宽恕之事。并且对于这一过失，他们在婚姻中得到的不是促使他们去犯这过失的东西，而是为此祈求宽恕的东西，只要他们不拒绝天主的仁慈，或是因在某些日子没有节制，使他们得以自由祈祷，并藉此节制的做法，如同藉着禁食，使他们的祈祷蒙悦纳；或是改变了自然的使用，变成了反乎本性的使用，后者发生在丈夫或妻子身上时更为可憎。\par

12. 因为，自然的使用，当它超越婚姻的契约时，即超越生育的必要性时，在妻子的情况下是可宽恕的，在妓女的情况下是该受惩罚的；而反乎本性的使用，发生在妓女身上是可憎的，但发生在妻子身上更为可憎。造物主的命令和受造物的秩序具有如此大的力量，以至于在允许我们使用的事物中，即使分寸过度，也远较在那些不允许的事物中，即使单独或罕见的过度，更容易容忍。因此，在允许的事情上，丈夫或妻子缺乏节制，是应该容忍的，以免情欲爆发到不允许的事情上。由此，一个不断与妻子亲近的人，比一个即使极少去行淫的人，所犯的罪要小得多。但是，当丈夫想要使用妻子的器官进行不被允许的用途时，如果妻子允许它发生在自己身上，她比发生在另一个女人身上更为可耻。因此，婚姻的装饰是生育的贞洁和履行肉身义务的信实；这是婚姻的本分，宗徒为它辩护，免受一切指责，他说：\BibleQuote{如果你娶了妻子，不算犯罪；如果童贞女嫁了人，也不算犯罪；} 以及：\BibleQuote{让她随心所欲，不算犯罪，只要她嫁了人。} 但是，出于上述原因，在要求对方履行配偶义务时过度超越节制，对已婚者而言被允许作为宽恕之事。\par

13. 因此，他所说的：\BibleQuote{没有丈夫的妇人，挂虑主的事，为能身心圣洁。}我们不应理解为一位贞洁的基督徒妻子在身体上不圣洁。因为对所有信徒都曾说过：\BibleQuote{你们不知道你们的身体是住在你们内的圣神的宫殿，你们从天主领受了这圣神吗？}因此，已婚者的身体也是圣洁的，只要他们对彼此和对天主保持忠信。而且，即使有一方是不信的，也不妨碍他们任何一方的这种圣洁，反而妻子的圣洁使不信的丈夫得益，丈夫的圣洁使不信的妻子得益，同一宗徒作证说：\BibleQuote{因为不信的丈夫因妻子成了圣洁的，不信的妻子因弟兄也成了圣洁的。}因此，那话是根据未婚者比已婚者有更大的圣洁而说的，相应地也有更大的赏报，因为一个是善，另一个是更大的善：同时她也只有这个挂虑，如何悦乐主。因为并非相信的女子，保持婚姻的贞洁，就不想如何悦乐主；但确切地说，她较少这样想，因为她挂虑世俗的事，如何悦乐丈夫。因为他正是要这样论她们：她们可以说，由于婚姻，被迫挂虑世俗的事，如何悦乐她们的丈夫。\par

14. 并且并非没有正当的理由，人们提出了一个疑问：他是对所有的已婚妇人说这话，还是对那样众多的人说，以至于几乎所有人都可以被认为是如此。因为就连他论未婚妇人的话：\BibleQuote{没有丈夫的妇人，挂虑主的事，为能身心圣洁，}也并非适用于所有未婚妇人：因为有些寡妇虽活着，却生活在放纵中。然而，就某种区分和她们（未婚与已婚）各自所具有的特征而言，正如那从婚姻中约束自己（即从允许的事中约束自己）却不约束自己于过犯（无论是奢侈、骄傲、好奇或多言）的人，应受极大的憎恨；同样，已婚妇人很少见到，她在婚姻生活的顺服中，除了如何悦乐天主之外别无他想，她不以辫发、黄金、珍珠和华服装饰自己，而是以善行，如那宣认虔诚的妇人所宜。诚然，宗徒伯多禄也在训导中描述了这样的婚姻。他说：\BibleQuote{同样，你们做妻子的，要顺服自己的丈夫；即使有人不信从道理，也可以因妻子的品行，不用言语，而被感化，因为看见了你们怀有敬畏的贞洁品行。你们的装饰不应是外表的辫发、佩戴金饰或华丽的衣裳，而应是那隐藏在心中的人，具有不朽的温柔与宁静精神，这在主前是宝贵的。因为从前那些盼望天主的圣妇们，也是这样装饰自己，顺服自己的丈夫：就如撒辣听从亚巴郎，称他为主；你们若行善，不怕任何惊吓，就成了她的女儿。同样，你们做丈夫的，要与妻子平安同居，尊重她们，因为她们是较软弱的器皿，与你们同享恩宠，免得你们的祈祷受阻碍。}难道这样的婚姻真的没有挂虑主的事，如何悦乐主吗？但它们非常罕见：谁能否认呢？而且，由于它们罕见，几乎所有这样的人，并不是为了成为这样的人而结合的，而是已经结合后才成为这样的人。\par

15. 如今我们时代的基督徒，既不受婚姻束缚，又能克制一切性交，看见正如经上所载，如今是\Quote{非拥抱之时，而是戒除拥抱之时}，难道不是宁可保持童贞或寡居的节欲，也不愿（既然没有对人类社会的义务所迫）忍受肉体之苦——没有这些痛苦，婚姻便无法存在（至于宗徒所省略的其他事，我们略而不提）吗？然而，当他们在欲望支配下结合之后，若此后克服了它，因为婚约一旦缔结便不可解除（正如当初不缔结是允许的），他们便成为婚姻形式所宣称的那样：或者通过彼此同意上升到更高的圣德境界，或者，若双方并非都是如此，那么那一个是如此的人就不会索求，而是给予所应得的，在一切事上保持贞洁而虔诚的和谐。但在那些时代，当我们的救恩奥迹尚隐于预像性的圣事中时，即使那些婚前已是如此的人，也因生养子女的职责而结婚，并非为情欲所胜，而是出于虔敬之心。倘若在当时就给予他们在新约启示中所给予的选择——主说：\Quote{谁能够领受，就领受吧}——那么，凡是仔细阅读过他们如何对待自己的妻子（当时甚至允许一人拥有多个妻子，其贞洁程度超过今日任何虽只有一个妻子却如我们所见宗徒以容许的方式所允许的那些人）的人，无人怀疑他们会欣然领受。因为他们是出于生养子女而拥有妻子，而非\Quote{像不认识天主的外邦人那样，陷于欲望的病症}。这是一件如此伟大的事，以至于如今许多人在一生中更容易戒除一切性交，而一旦结婚，却难以节制到仅为子女而行的程度。我们有许多弟兄和姊妹，无论男女，都是节欲之人——或是曾经历过婚姻的，或是完全远离一切此类关系的——他们不计其数。然而，在我们与他们的日常交谈中，我们曾听见过谁——无论是已婚者还是曾结婚者——告诉我们，他从未与妻子行房事除非是怀着受孕的希望？因此，宗徒所命令已婚者的事，才是婚姻所固有的；而他们所宽恕的事，或阻碍祈祷的事，婚姻并非强迫，而是容忍。\par

16. 因此，倘若偶然——（我不知道这是否可能发生，而且我倾向于认为不可能发生；但倘若偶然）——一个男人暂娶一名姘妇，只是为了从这种交合中求得子嗣，那么这种结合也不能优于那些妇女的婚姻，因为后者所做的事是应受宽恕的。因为我们必须考虑什么是婚姻固有的，而不是考虑那些嫁人却以不够节制的方式使用婚姻的妇女所做的事。因为，即使有人不正当地、错误地占有土地，却从土地的出产中大量施舍，他也不能因此证明抢掠是正当的；同样，若有人因贪婪而霸占继承得来的或正当获得的田产，我们也不能因此指责民法规则，正是这规则使他成为合法的所有者。暴君反叛的不义也不应受赞美，假如暴君以王者的仁慈对待臣民；同样，王权的秩序也不应受指责，假如国王以暴君的残酷施政。因为，愿意善用不正当的权力是一回事，而不正当地运用正当的权力是另一回事。因此，临时娶来的姘妇，即使是为了子嗣，也不能使她们的姘居合法化；已婚妇女若与丈夫放荡生活，也不能将任何责任归咎于婚姻秩序。\par

17. 婚姻可以由起初不正当结合的人缔结，随后以诚实的决议使之成立，这是显而易见的。但在我们天主的城中，一旦缔结婚姻——甚至从一男一女最初的结合起，婚姻就带有某种圣事的性质——除了其中一人的死亡，绝不能解除。因为婚姻的纽带依然存在，即使由于明显的不孕而没有子女——子女原是婚姻缔结的目的；这样，当已婚者知道他们不会有孩子时，即便为了子女的缘故，他们也不能合法地分离而与他人结合。如果他们如此行，便同那些与之结合的人犯了奸淫，但自己仍然保持夫妻关系。显然，在古时的先祖中，在妻子的同意下另娶一个女人，从她那里为两人共同生育儿子——由一人的性交和种子，但由另一人的权利和权力——是合法的。至于现今是否也合法，我不愿轻易断言。因为如今不再有生养子女的必要，如同那时一样；那时即使妻子生育了子女，为了更多后裔的繁衍，也允许再娶妻妾，而现今这当然是不合法的。因为时代的差异使得适时性对于正义以及做与不做某事有如此大的影响，以至于现在一个男人若不娶妻（除非他不能节制），反而做得更好。但那时他们甚至娶好几个妻子也毫无可责之处，即使那些更能克制的人也是如此，若不是因为那时的虔诚另有要求。因为，正如那位智慧而正义的人，如今渴望离世与基督同在，并更以这最好之事为乐，如今不是出于生活在此地的欲望，而是出于对他人有益的职责而进食，以便留在肉身中；同样，以婚姻的权利与女性交合，对那时的圣人而言是职责，而非情欲。\par

18. 因为食物是为保存个人，性交是为保存种族；两者都离不开肉身的快感；但这快感若加以调节，并由节制的约束回归本性之用，就不能称为情欲。正如非法食物在维持生命时是淫乱或奸淫的性交在寻求子嗣时的情况；以及非法食物在肚腹和喉咙的奢侈中，非法性交在不寻求子嗣的情欲中；以及在合法食物上某些人的过度食欲，相当于在夫妻之间那种可宽恕的性交。因此，与其吃祭偶像之物，宁可饿死；同样，与其从非法性交寻求子嗣，宁可无子而死。但无论人从何而生，若他们不效法父母的恶习，并正确敬拜天主，他们将诚实而安全。因为人的种子，无论出自何种人，都是天主的创造；滥用种子的人将遭恶报，但种子本身决不会成为邪恶。正如奸淫者的好儿子不能为奸淫辩护，同样已婚者的坏儿子也不能归咎于婚姻。因此，正如新约时期的先祖们出于保存的责任取用食物，虽然他们带着肉身的自然快感，但决无法与那些吃祭偶像之物者的快感相比，也不可与那些虽然食物合法却过量取用者的快感相比；同样，旧约时期的先祖们出于保存的责任使用性交；而他们自然的快感，决未放纵于不合理和非法的情欲，不能与淫行的卑劣或已婚者的无节制相比。的确，通过同一股爱德之脉，如今按精神，那时按肉身，为那母亲耶路撒冷之故生育儿子是一种责任；除了时代的不同，别无其他使先祖们的作为不同。但这样，甚至先知们，虽不按肉身生活，也必须按肉身交合；同样，宗徒们，虽不按肉身生活，也必须按肉身进食。\par

19. 因此，现今许多妇女，对她们说\BibleQuote{若她们不能自制，就让他们结婚}，不能与当时的圣妇相比，即使她们也结婚。婚姻本身在万民中都是为了生育子女，无论子女后来成为怎样的人，婚姻原是为这目的设立的，使他们能以适当而有序的方式出生。但那些不能自制的人，可以说是藉着一级诚实的阶梯上升到婚姻；而那些肯定能自制的人，如果那时代的目的允许这样，则在某种程度上下到婚姻，藉着一级虔诚的阶梯。因此，虽然两者的婚姻，就其作为婚姻而言，因为是为了生育，是同样好的，但这些已婚的人不能与那些已婚的人相比。因为这些人拥有，因婚姻的诚实而允许他们的，虽不属于婚姻；即超出生育必要性的部分，而那些人没有。但若现在碰巧有人，在婚姻中除婚姻设立的目的外别无所求、别无所欲，也不能与那些人同等。因为在这些人的心中，对儿子的渴望是肉性的，但在那些人的心中是精神的，因为它适合于那时代的圣事。的确，如今没有人在虔诚上达到成全，会按肉性的意义寻求生子；但那时，按肉性的意义生子本身就是虔诚的工作，因为那个民族的生育充满了未来事物的预告，并属于先知的安排。\par

20. 因此，正如允许一个男人有几个妻子，却不允许一个女人有几个丈夫，即使是为了子嗣，以防妇女能生育而男人不能生育。因为自然有一种隐秘的法则：为首的事物喜欢单一，但附属的事物被置于其下，不仅是一对一，而且如果自然或社会的体系允许，甚至多个附属在一个之下，也不失其美。因为没有一个奴隶有多个主人，如同多个奴隶有一个主人那样。所以我们没有读到任何一位圣妇服侍两个或更多活着的丈夫；但我们读到许多妇女服侍一个丈夫，当那民族的社会状况允许，且那时代的目的劝导这样时：因为这并不违反婚姻的本质。因为多个妇女可以从一个男人怀孕；但一个妇女不能从多个男人怀孕，（这便是为首之物的能力：）正如许多灵魂理应臣服于唯一的天主。因此，灵魂的唯一真天主只有一位；但一个灵魂可以通过许多假神行淫，却无法结出果实。\par

21. 然而，既然从许多灵魂中，将来会有一座城，由那些同有一心一意归于天主的人组成；我们未来完美的合一，是在异乡漂泊之后才能实现的，那时，所有人的思想既不会彼此隐藏，也不会在任何事上相互对立；正因如此，我们时代的婚姻圣事已被缩减为一夫一妻，以致除了一妇之夫外，不得祝圣任何人作教会的管家。那些认为连作为慕道者或外教人时曾有第二个妻子的人也不应被祝圣的人，对此理解得更为深刻。因为这是关于圣事，而非关于罪。因为在圣洗中，一切罪都被除去。但那位说：\BibleQuote{你若娶妻，并没有犯罪；若处女出嫁，也没有犯罪}以及\BibleQuote{她若愿意，就由她，她若出嫁，也没有犯罪}的，已清楚表明婚姻不是罪。但因圣事的圣洁，正如一个女子，即使作为慕道者时遭受过强暴，在领洗后也不能被祝圣入天主的贞女行列；同样，认为那些超过一个妻子的人，并非犯了什么罪，而是失去了某种圣事的规范，这种规范并非为善生的功德所必需，而是为教会祝圣的印记所必需，这也并不荒谬。因此，正如古圣祖们多有妻子，预表我们将来从万民中归于同一丈夫基督的教会；同样，我们的大司祭（一妇之夫）预表从万民中归于同一丈夫基督的合一；这合一将在祂暴露黑暗中的隐密、显明人心的思念时得以圆满，那时各人要从天主那里得着称赞。但现在，即使在那些将来要合而为一、归于一体的人中间，也有明显的或隐藏的不和，虽然爱德仍安全无恙；这些不和将来肯定不复存在。因此，正如那个时代的多妻婚姻圣事预表将来在地极万民中要归于天主的众多子民，同样我们这个时代的一夫一妻婚姻圣事预表我们所有人将来在天上之城归于天主的合一。因此，如同事奉两个或更多，或从活着的丈夫转嫁另一人，那时不允许，现在也不允许，将来也永不允许。事实上，背弃独一真天主，去行另一位的奸淫迷信，永远是恶事。因此，我们的圣人们并非为了子女更多而行那罗马人加图据说所行的事：在有生之年休妻，以便甚至用儿子充满别人的家。因为在一妇之婚中，圣事的神圣比生育的果实更有价值。\par

22. 因此，如果甚至连那些仅为生育而结合的人——婚姻正是为此目的而建立——也不能与圣祖们相比，因为圣祖们寻求自己儿子是在一种远非世俗的方式中；例如亚巴郎，当受命去杀自己的儿子时，无畏而虔诚，没有放过自己的独生子，那个他从极大绝望中得来的儿子，只是在那位命令他举手者的禁止下才放下了手；那么，我们就要思考，至少我们中间节欲的人是否能与那些已婚的圣祖们相比；除非现在这些人应该被优先于那些我们尚未找到可与之相比的人。因为他们的婚姻中有更大的益处，超过婚姻本身的益处：无疑，节欲的益处应优先于这些：因为他们并非出于像今天的人那样的义务从婚姻中寻求儿子——那种义务源于某种要求延续以对抗死亡的必死本性的意识。而否认这是善的人，不认识天主，一切善物的创造者，从天上到地上，从不朽到有死。但连野兽也不完全没有这种生育的意识，特别是鸟类，它们筑巢的辛劳立即向我们显现，还有某种类似婚姻的东西，为了共同生育和养育。但那些人，以远为圣洁的心思，超越了必死本性的这种情感——这种情感在其同类中的贞洁，加上对天主的敬拜，如同一些人理解的那样，被描述为结出先三十倍的果实；他们为基督的缘故寻求婚姻中的儿子，为要将祂按血统的种族从万民中分别出来：正如天主喜悦这样安排，使这事超越其他一切有利于预言祂，因为预言了祂将来按血肉应出自哪个种族和哪个民族。因此，这远胜于我们中间信徒的贞洁婚姻，亚巴郎在他的大腿下知道这一点，他叫仆人把手放在那里，好起誓关于他将要娶妻的妻子。因为把手放在一个人的大腿下，并指着天上的天主起誓，若非指明那从大腿起源的血肉中，天上的天主将要来临，还有什么别的意思呢？因此，婚姻是善的，其中已婚的人越是以更大的贞洁和信德敬畏天主，他们就越善，特别是如果他们按血肉所渴望的儿子，也按精神来教养。\par

23. 法律命人在与妻子同房后也要取洁，这并不表明那是罪，除非那是因宽恕而被允许的，而这宽恕若过度也会阻碍祈祷。然而，由于法律将许多事物设立在圣事和未来之事的预象中，种子的某种仿佛质料无形状态\OriginalTerm{质料无形状态}{material formless state}，在获得形式后将产生人的身体，被用来象征一种无形无教化的生命；人既有必要藉着学习和教导的形式与教化而得洁净，因此作为此事的预象，法律在射精后命令了取洁\OriginalTerm{取洁}{purification}。因为即使在睡眠中，这事也不是因罪发生的。然而那里也命令了取洁。或者，若有人认为这也是罪，以为它只是由某种这类情欲而来（这无疑是错误的）；那么，妇女的普通月经也是罪吗？然而，旧法律却命令她们藉赎罪\OriginalTerm{赎罪}{expiation}而取洁；这没有别的原因，只由于那质料无形状态本身——在怀孕发生之后，它仿佛被添加来建造身体，而因此当它无形地流出时，法律愿意藉此象征一个没有纪律形式的灵魂，不体面地流动和散漫。而它应当获得形式，当法律命令这样的身体流动应当被取洁时，便表明了这一点。最后，那么，死亡也是罪吗？或埋葬死者，不也是人道的善工吗？然而，即使在这种情况下也命令了取洁；因为死尸，生命离弃了它，并不是罪，却象征着被义德离弃的灵魂之罪。\par

24. 婚姻，我说，是善的，并且可以用正确的理性抵御一切诽谤。但是论到圣祖们的婚姻，我所探究的不是哪种婚姻，而是哪种节德\OriginalTerm{节德}{continence}处于同一水平；或者更恰当地说，不是拿婚姻与婚姻比较（因为赐予必死人性的恩赐在所有情形下是相等的），而是那些使用婚姻的人——因为我发现他们与其他以截然不同的精神使用婚姻的人无法相比——我们必须探究，节德之人承认自己可以与哪些已婚之人相比。除非，亚巴郎\Person{亚巴郎}或许不能为天国之故禁戒婚姻吗？他为了天国之故，能无所畏惧地祭献他唯一获得后裔的抵押品——因了这后裔，婚姻才是宝贵的！\par

25. 诚然，节德不是身体的德行，而是灵魂的德行。但灵魂的德行有时在行动中显示，有时隐藏在习德\OriginalTerm{习德}{habit}中，就如殉道的德行藉忍受苦难而闪耀显现；但有多少人拥有同样的心灵德行，却缺少考验，使得内心在天主面前的事物也能显露在人面前，并且不是在人面前开始存在，而只是变得为人所知？因为约伯早已有忍耐，天主知道并为他作证；但藉着试炼的考验，忍耐才为人所知：内在隐藏的并非被产生，而是藉着从外面加于他的事物而显出。提摩太也肯定有戒酒的德行，保禄并未从他那里夺去这德行，而是劝他\BibleQuote{为了他的胃和屡次患病}而适量用点酒；否则他便教他一堂致命的课：为了身体的健康而损失灵魂的德行。但是因为他的劝告可以在那德行的安全下实行，饮用的益处便如此自由地留给身体，而节德的习德继续存在于灵魂中。正是这习德本身，使人当需要时做某事；但当不做时，也能够做，只是没有需要。在关于性交的节德方面，那些被说\BibleQuote{他们若不能节制，就让他们结婚}的人没有这习德。但那些被说\BibleQuote{谁能领会，就让他领会}的人有这习德。因此，完美的灵魂藉着这节德的习德，如此使用那些为别事所需的现世财富，以致不被它们束缚，并且由此有力量在不需要时也不使用它们。除非谁也有力量不使用它们，否则无人能善用它们。的确，许多人更容易实践禁欲而不使用，却难以实践节制\OriginalTerm{节制}{temperance}而善用。但除了能节德地不使用它们的人，无人能明智地使用它们。从这习德，保禄也说\BibleQuote{我知道怎样处卑贱，也知道怎样处丰富}。诚然，处卑贱是任何人的事；但知道怎样处卑贱是伟人的事。同样，处丰富，谁不能呢？但知道怎样处丰富，只有那些不被丰富腐化的人才能。\par

26. 但是，为了更清楚地理解，在\OriginalTerm{习性}{habitus}中如何能有德行，即使它不在\OriginalTerm{行动}{opus}中，我要举一个例子，关于这一点，没有一个天主教基督徒会怀疑。因为我们的主\Person{耶稣基督}在真实肉身上曾饥渴、饮食，凡信福音的人无人怀疑。那么，在祂里面难道没有像在\Person{洗者若翰}里面那样大的饮食\OriginalTerm{节德}{virtus continentiae}吗？因为若翰来，不吃不喝，人们却说：他附了魔；人子来，也吃也喝，人们却说：\BibleQuote{看，一个贪吃好酒的人，税吏和罪人的朋友。}难道不也正是用另一种方式，就性交而言，有人对我们的祖宗、祂的家人说：\Quote{看，好色不洁的人，贪爱女人和淫荡的人？}然而在祂身上，这并不真实，尽管祂确实不像若翰那样禁食禁饮，因为祂自己最清楚最真实地说：\BibleQuote{若翰来了，不吃不喝；人子来了，也吃也喝：}同样，在这些祖宗身上也并不真实；虽然现在来了基督的宗徒，不婚娶，不生育，以致外教人说他是魔法师；但那时却来了基督的先知，娶妻生子，以致\OriginalTerm{摩尼教徒}{Manichaei}说他是贪恋女色的人：\BibleQuote{而智慧，祂说，由她的子女得了证明。}主在那里说了关于若翰和祂自己之后，接着说：\BibleQuote{但智慧，祂说，由她的子女得了证明。}这些人看到，\OriginalTerm{节德}{virtus continentiae}应当在灵魂的习性中存在，但要根据事物和时机的机会在行动中显示出来；就如圣殉道者们的忍耐德行在行动中显现；而同样圣洁的其他人则在习性中。因此，就如受苦的\Person{伯多禄}和未受苦的\Person{若望宗徒}的忍耐功绩并非不平等；同样，未尝试婚姻的\Person{洗者若翰}和生子的\Person{亚巴郎}的节德功绩也并非不平等。因为若翰的独身和亚巴郎的婚姻都按照时代分配，为基督服役；但若翰在行动上也有节德，而亚巴郎仅在习性上。\par

27. 因此，在那时期，当法律在列祖的日子之后，也诅咒凡不在以色列中立嗣的人，即使那些能够这么做的人，也不付诸行动，但仍然拥有它。但从时期一满，就应当说：\BibleQuote{谁能领受，就领受吧。}从那时起直到现在，并从此直到终末，凡拥有的人就在行动；凡不愿行动的人，就不要谎称自己拥有。藉此，那些以恶言败坏善行的人，以空虚狡猾的诡计，向一个持守\OriginalTerm{节德}{continentia}、拒绝婚姻的基督徒说：那么，你比亚巴郎更好吗？但他听到时，不要惊惶；也不要胆敢说：\Quote{更好，}也不要脱离他的目标：因为前者他说得不真实，后者他做得不对。但让他说：我确实不比亚巴郎好，但\OriginalTerm{独身贞洁}{castitas caelibatus}比\OriginalTerm{婚姻贞洁}{castitas coniugalis}更好；亚巴郎在习性上拥有两者，在行动上使用了其中之一。因为他在婚姻状态中贞洁地生活；但在他能力之内，他可以无婚姻而贞洁，只是那时不适宜。而我更容易不使用亚巴郎所使用的婚姻，而不能像亚巴郎那样使用婚姻：因此，我比那些因心神\OriginalTerm{失节}{incontinentia}而不能做我所做的事的人好；而不是比那些因时代不同而没有做我所做的事的人好。因为我现在所做的，如果那时必须做，他们会做得更好；但他们所做的，我不应该那样做，即使现在必须做。或者，如果他感觉并知道自己是这样的人，在保持并继续心智习性中的节德的情况下，如果因某种宗教义务而下降到婚姻的使用中，他将会是像\Person{亚巴郎}那样的丈夫和父亲；那么让他勇敢地回答那个刁难的问题者，并说：我确实不比亚巴郎好，只是在这种节德方面，他并不缺乏，尽管没有显现：但我就是这样的人，没有别的，只是所做的不同。让他这样说清楚：因为即使他愿意夸耀，他也不愚蠢，因为他讲真话。但如果他有所保留，以免有人因他所见或所闻而对他有更高的看法；让他从自身除去问题的症结，不是针对人，而是针对事情本身来回答，并说：凡有如此大能力的人，就是像亚巴郎这样的人。但可能发生的是，节德在他心中较小，他未使用亚巴郎所使用的婚姻：然而，它比那因不能更大而持守婚姻贞洁的人心中的节德更大。同样，那思想主的事、好使她身体和灵魂都圣洁的未嫁女子，当她听到那无耻的问题者说：那么，你比\Person{撒辣}更好吗？她应回答：我更好，但比那些缺乏节德的人好，我相信撒辣并不缺乏：她因此连同这种德行做了适合那时代的事，而我得以免除，使我的身体也能显现她心中所保持的。\par

28. 因此，如果我们比较事物本身，我们绝不能怀疑独身贞洁比婚姻贞洁更好，虽然两者都是善的；但当我们比较人时，拥有更大善的人就更好。此外，同一种类中拥有更大的人，也拥有那较小的；但只拥有较小的人，肯定不拥有那更大的。因为在六十中也包含三十，但在三十中却不包含六十。但不从他所拥有的出发去行动，乃在于职责的分配，不在于德行的缺乏；因为那找不到值得他慈悲援助的可怜人的人，也没有缺乏\OriginalTerm{慈悲的善}{bonum misericordiae}。\par

29. 还有这一点：人不能单单以某一种善来彼此比较。因为可能发生这样的事：一个人没有另一个人所有的，却拥有另一件应被视为更有价值的东西。\OriginalTerm{服从}{obedience}之善胜于\OriginalTerm{节欲}{continence}之善。因为婚姻从未被我们圣经的权威所谴责，而\OriginalTerm{不服从}{disobedience}却从未被开脱。因此，如果有一个愿意保持\OriginalTerm{童贞}{virginity}的童贞女，但却是\OriginalTerm{不服从}{disobedience}的，而有一个已婚妇人，她不能保持童贞，却是服从的，那么我们将说哪一个更好呢？是那个（童贞女）不如她若保持童贞时值得称赞呢，还是那个（已婚妇人）即使她是童贞也值得责备呢？同样，如果你将一个醉酒的童贞女与一个清醒的已婚妇人比较，谁能怀疑会做出同样的判断呢？的确，婚姻与童贞是两种善，其中一种更大；但清醒与醉酒，正如服从与倔强，一个是善，另一个是恶。然而，拥有较小程度的所有善，胜过拥有巨大的善伴随巨大的恶：因为在身体之善中也是如此，拥有\Person{匝凯}的身材且健康，胜过拥有\Person{哥肋雅}的身材却发烧。\par

30. 问题的关键显然不是：一个各方面都不服从的童贞女是否可与一个服从的已婚妇人相比较，而是一个较不服从的与一个较服从的比较：因为婚姻也有其贞洁，因此是一种善，但小于童贞。因此，如果一个人，在服从之善上越小，在贞洁之善上越大，与另一个人比较，那么谁应被优先选择，那首先将贞洁本身与服从本身加以比较的人会判断出，他看到服从在某种意义上是一切德行之母。因此，基于这个理由，可以有不带童贞的服从，因为童贞是劝谕，而非\OriginalTerm{诫命}{precepts}。但我所称的服从，是遵行诫命的服从。因此，可以有不带童贞的诫命服从，但不能没有贞洁。因为\OriginalTerm{贞洁}{chastity}关乎不犯奸淫、不犯邪淫、不被任何非法交合所玷污：凡不遵守这些的，就是违反天主的诫命，并因此被逐出服从的德行。但可以有不服从而童贞，这是因为一个女人可能接受了童贞的劝谕并守住了童贞，却轻视诫命：正如我们认识的许多神圣的童贞女，多言、好奇、酗酒、好争、贪婪、骄傲：这一切都违反诫命，并因\OriginalTerm{不服从}{disobedience}之罪而杀死人，正如\Person{厄娃}本人一样。因此，不仅服从者应优先于不服从者，而且更服从的已婚妇人应优先于较不服从的童贞女。\par

31. 由于这种\OriginalTerm{服从}{obedience}，那位并非没有妻子的父亲（\Person{亚巴郎}）已准备失去独生子，并且是由他亲手杀死。因为我并非无端地称他为独生子，关于他，他听见主说：\BibleQuote{因为借着依撒格，你的后裔才得传名}。因此，如果他被命令连妻子也要失去，他会多么迅速地听从呢？因此，我们常常有理由认为，有些男女，在\OriginalTerm{节欲}{continence}、远离一切性交后，却因如此热切地抓住不使用被允许之事，而在遵守\OriginalTerm{诫命}{precepts}方面疏忽。那么，谁怀疑我们不应把我们的同时代人——虽然远离一切交合，但在服从之德上较低——与那些生养儿子的圣父和圣母的美善相比较呢？即使那些人在心思习惯上也有所欠缺，而这在后者的事迹中是明显的。因此，让这些人跟随羔羊，男孩子唱着新歌，正如\WorkTitle{默示录}所记载的：\BibleQuote{这些人没有与女人玷污自己}。没有别的原因，只是因为他们保持了\OriginalTerm{童贞}{virginity}。也不要因此以为自己比那些最初使用婚姻的圣父们更好，可以说是以婚姻的方式。确实，婚姻的使用是这样的：如果其中通过肉体的交合发生了任何超出生育需要的事，虽然是可以宽恕的，但仍有\OriginalTerm{污秽}{pollution}。如果这种过度没有造成任何污秽，那么宽恕又赎了什么？如果那些跟随羔羊的男孩子没有保持童贞，他们免于这种污秽就是奇怪的。\par

32. 因此，婚姻之善在所有民族和所有人中，皆在于生育后代的机会和贞洁的忠贞；然而，就属于天主子民者而言，也在于圣事的圣洁性，因这圣洁性，离开丈夫的妇人，即使她被休弃，在其夫在世时，再嫁给另一个人是不合法的，甚至为了生育子女也不可。而且，尽管这是婚姻的唯一原因，但即使那导致婚姻的事情没有发生，婚姻的纽带也不会解开，除非夫妻一方死亡。同样，就像为了形成会众而举行圣职授任，即使没有会众形成，但在被授任者身上仍保留着圣秩圣事；并且，如果因任何过失而有人被撤职，他也不会失去那一次而永远被置于他身上的主的圣事，尽管这使他被定罪。因此，婚姻是为了生育子女，宗徒为此作证说：\Quote{我愿意}\newline{}他说：\Quote{年轻的妇人结婚。}\newline{}并且，好像有人问他：为了什么目的？他立刻补充说：\Quote{生儿育女，做一家之母。}\newline{}但关于贞洁的忠贞，有这样的话：\Quote{妻子对自己的身体没有主权，但在丈夫有；同样，丈夫对自己的身体也没有主权，但在妻子有。}\newline{}至于圣事的圣洁性，有这样的话：\Quote{妻子不可离开丈夫；若是离开了，就当不再嫁人，或是与丈夫和好；并且丈夫也不可休妻。}\newline{}所有这些都是善，婚姻因之而成为善：即子女、忠贞、圣事。不过如今，不寻求肉体的后代，并借此维持某种脱离一切此类工作的永久自由，且在精神上归属于唯一的丈夫基督，这无疑是更好、更神圣的；条件是，人们如此使用这份自由，正如经上所记：默思主的事，怎样悦乐主；也就是说，节德要时刻注意，勿使服从在任何事情上有所欠缺。而这一德行，作为根本之德，及（通常所谓的）母德，且显然是普遍的，古代圣父们实际上已加以实践；但他们在心灵习性中拥有节德。无疑，他们借着那使他们成为正义和圣洁的服从，并随时准备行一切善事，即使他们被命令完全禁绝所有房事，也会执行。因为，若他们在天主的命令或劝勉下不进行房事，那该多么容易啊——他们为了服从，能杀死那唯一通过房事而生的孩子，而房事正是为了这孩子的生育。\par

33. 既然如此，对于异端者，无论是摩尼教徒，还是其他任何人，就旧约圣父们拥有多位妻子之事对他们提出的指控，我们已经做出了足够甚至过多的答复。他们认为这是证明圣父们不节制的证据。也就是说，只要他们认识到：那既不违犯自然（因为他们与那些女人相交并非出于放纵，而是为了生育子女），也不违犯习俗（因为那些事在那时是常见的），也不违犯诫命（因为没有法律禁止）的行为，就不是罪。而那些用非法方式与女人相交的人，或是在那些圣经中被天主的审判所定罪，或是经籍指出他们当受谴责和避之不及，而不是加以赞许或效仿。\par

34. 但那些我们有妻子的弟兄们，我们尽力劝诫他们，切勿以自己的软弱来判断那些圣父们，正如宗徒所说，他们拿自己与自己比较；因此，他们不明白灵魂在侍奉正义对抗私欲方面有多大的力量，它不会顺从这类肉体的冲动，也不会容许它们滑向或推进到超出生育子女必要性的房事，只要遵循自然的秩序、习俗的惯例、法律的命令。确实，人们对那些圣父们有这样的猜疑，是因为他们自己或因不节制而选择了婚姻，或对自己的妻子放纵无度。然而，那些节德的人，无论是男人（在妻子去世后）、女人（在丈夫去世后），或是双方经相互同意向天主发了节德誓愿的人，要知道，他们理应得到比婚姻的贞洁所要求的更大的赏报；但至于那些圣父们的婚姻，他们是以预言的模式结合，在房事中只寻求子女，在子女中只寻求那将要来临的、将来在肉身中来到的基督，他们不仅不应因自己的生活方式而鄙视这些婚姻，反而要毫无疑问地将其置于自己的生活方式之上。\par

35. 童男与童贞女，向天主献上实际的贞洁，我们首先告诫他们，使他们知道，他们必须以极大的谦卑来守护现世的生命，因为他们所发之愿越是属天，就越当如此。诚然经上记载：\Quote{你无论多么伟大，在一切事上都要自谦自卑。}因此，我们有责任论说他们的伟大，他们则有责任思念极大的谦卑。所以，除了某些结过婚的圣祖与圣妇以外——这些人虽然尚未结婚，并不比他们更好，因为即使他们结了婚，也无法和那些圣祖圣妇相等——他们不要怀疑，自己超越了同时代所有其他的人，无论是已婚者，还是在尝试婚姻生活后修习节欲者；但并非像\Person{亚纳}超越\Person{苏撒纳}那样，而是像\Person{玛利亚}超越她们二者那样。我所说的是关于肉身的神圣贞洁本身；因为谁不知道\Person{玛利亚}还有其他功绩呢？因此，他们须将相称的行为加于如此崇高的目标上，这样他们才能确信那非凡的赏报；他们确知，对于他们自己以及所有信者、基督所爱所选的肢体，有许多将从东方和西方而来，虽然按各自的功绩闪耀着彼此不同的荣光，却共同领受这份伟大恩赐，即与\Person{亚巴郎}、\Person{依撒格}、\Person{雅各伯}一同坐在天主的国里；而这些人，并非为了世俗，而是为了基督才成为丈夫，为了基督才成为父亲。\par

\SourceRecord{Translated by C.L. Cornish.；Nicene and Post-Nicene Fathers, First Series；Vol. 3.；Edited by Philip Schaff.；Buffalo, NY: Christian Literature Publishing Co.,；1887.；<http://www.newadvent.org/fathers/1309.htm>.}
